\documentclass[11pt]{article}
\usepackage[utf8]{inputenc}
\usepackage[T1]{fontenc}
\usepackage[margin=2.5cm]{geometry}
\usepackage{graphicx}
\usepackage{float}
\usepackage{hyperref}
\usepackage{array}

\title{Identificación de subgrupos clínicos de COVID-19 mediante expresión génica}
\author{Rivera Morales David \\ Lopez Bernal Yeimi Lizet \\ Hidalgo Carrillo Amir Gilberto \\ Badager Estrada Aaron Omar}
\date{Genómica Computacional -- Semestre 2026-II\\GSE157103, GEO/NCBI}

\begin{document}
\maketitle

\section*{Resumen}
Este proyecto analiza GSE157103 para revisar si la expresión génica en sangre
permite formar grupos parecidos a UCI/no UCI en pacientes con COVID-19. Se
usaron matrices procesadas de conteos esperados y TPM, se filtraron genes con
baja señal y se trabajó con $\log_2(TPM + 1)$. Después se compararon k-means,
agrupamiento jerárquico y DBSCAN. El mejor resultado fue agrupamiento
jerárquico con ward, k=2, con
ARI=0.263, NMI=0.207,
silueta=0.176 y exactitud binaria aproximada de
0.76. La señal no es perfecta, pero sí es consistente:
un grupo concentra más pacientes UCI y el otro más pacientes no UCI.

\section{Introducción}
COVID-19 no se presenta igual en todos los pacientes. Algunos cursan la
enfermedad con síntomas moderados y otros requieren cuidados intensivos. Esa
diferencia clínica puede reflejar cambios moleculares en sangre, sobre todo en
genes relacionados con inflamación y respuesta inmune. La pregunta de este
proyecto fue directa: si no se le dan al algoritmo las etiquetas UCI/no UCI, ¿los
perfiles de expresión forman grupos parecidos a esas categorías?

\section{Datos y metodología}
Se usaron tres archivos de GSE157103: conteos esperados, TPM y metadatos de GEO.
Después de alinear expresión y metadatos quedaron 126 muestras:
100 COVID-19 y 26 controles no COVID-19. Dentro de las
muestras COVID-19 hay 50 pacientes UCI y
50 no UCI. La comparación principal se hizo
solo con pacientes COVID-19.

Se retuvieron genes con TPM al menos 1 y conteos esperados al menos 10 en un
mínimo de 5 muestras COVID-19 o 10\% de las muestras COVID-19. Después del
filtrado quedaron 11924 genes. La expresión se transformó
como $\log_2(TPM + 1)$ y se seleccionaron los 3000 genes de
mayor varianza.

Se aplicaron PCA y UMAP para visualización. Para agrupamiento se compararon
k-means, agrupamiento jerárquico y DBSCAN. La concordancia con UCI/no UCI se
evaluó con ARI, NMI, silueta y exactitud binaria cuando el método produjo dos
grupos.

\subsection*{Procesamiento aplicado}
\begin{enumerate}
\item Se filtraron genes con baja expresión.
\item Se aplicó la transformación $\log_2(TPM + 1)$.
\item Se seleccionaron los 3000 genes más variables.
\item Se aplicó PCA y UMAP.
\item Se probaron k-means, agrupamiento jerárquico y DBSCAN.
\item UCI/no UCI se usó solo para evaluar los grupos, no para construirlos.
\end{enumerate}

\section{Resultados}
El mejor método por ARI fue el agrupamiento jerárquico con
ward, k=2. El resultado fue:

\begin{center}
\begin{tabular}{l r}
\hline
Métrica & Valor \\
\hline
ARI & 0.263 \\
NMI & 0.207 \\
Silueta & 0.176 \\
Exactitud binaria & 0.76 \\
\hline
\end{tabular}
\end{center}

\subsection*{Comparación por método}
\begin{center}
\small
\begin{tabular}{l r l r r r r}
\hline
Método & Corridas & Mejor configuración & Grupos & ARI & Silueta & Ruido \\
\hline
DBSCAN & 44 & min\_samples=3, eps=22.71 & 6 & 0.060 & -0.055 & 0.57 \\
jerárquico & 15 & ward, k=2 & 2 & 0.263 & 0.176 & 0.00 \\
k-means & 5 & k=3 & 3 & 0.226 & 0.229 & 0.00 \\
\hline
\end{tabular}
\end{center}

DBSCAN se evaluó en 44 combinaciones de parámetros. Su mejor ARI fue 0.060 con min\_samples=3, eps=22.71 y fracción de ruido 0.57. Esto deja claro que DBSCAN sí se probó, pero no
fue el método más útil para esta pregunta.

\subsection*{Tabla grupo vs. UCI}
\begin{center}
\begin{tabular}{r r r}
\hline
Grupo & UCI & No UCI \\
\hline
0 & 40 & 14 \\
1 & 10 & 36 \\
\hline
\end{tabular}
\end{center}

\begin{itemize}
\item \textbf{Grupo 0}: 40/54 UCI (74.1\%) y 14/54 no UCI (25.9\%).
\item \textbf{Grupo 1}: 10/46 UCI (21.7\%) y 36/46 no UCI (78.3\%).
\end{itemize}

Observación: el grupo 0 concentra más pacientes UCI y el grupo 1 concentra más
pacientes no UCI. La separación no es perfecta, pero sí hay una dirección clara.

\subsection*{Genes marcadores principales}
\begin{itemize}
\item \textbf{Grupo 0}: IL1R2, C3orf86, UPP1, OLAH, GRB10, ROPN1L, GMFG, LRG1, FLOT1, RSPH14
\item \textbf{Grupo 1}: CD4, TMEM109, AARS2, ST3GAL5, TP53, CHST14, DCP1B, PATZ1, EID2, RSAD1
\end{itemize}

\subsection*{Enriquecimiento funcional}
\begin{itemize}
\item \textbf{Grupo 0}: degranulación de neutrófilos; respuesta a estímulos; respuesta inflamatoria; respuesta de defensa; sistema inmune innato
\item \textbf{Grupo 1}: diferenciación de células Th17; diferenciación de células Th1 y Th2; metabolismo de RNA no codificante; ensamblaje de snRNP; infección por virus de leucemia de células T humanas tipo 1
\end{itemize}

Observación: el grupo 0 tiene una señal más inflamatoria. Genes como
\textit{IL1R2}, \textit{UPP1}, \textit{OLAH} y \textit{LRG1} son
compatibles con activación inmune innata y cambios mieloides. El grupo 1 incluye
\textit{CD4} y términos relacionados con células T, por lo que parece menos
dominado por inflamación innata.

\section{Figuras principales}
\begin{figure}[H]
\centering
\includegraphics[width=0.95\linewidth]{\detokenize{../results/figures/qc_library_detected_genes.png}}
\caption{Control de calidad por tipo de muestra y estado UCI.}
\end{figure}

\begin{figure}[H]
\centering
\includegraphics[width=0.95\linewidth]{\detokenize{../results/figures/pca_covid_icu_clusters.png}}
\caption{PCA de muestras COVID-19 coloreado por estado UCI y grupo seleccionado.}
\end{figure}

\begin{figure}[H]
\centering
\includegraphics[width=0.85\linewidth]{\detokenize{../results/figures/clustering_ari_top_methods.png}}
\caption{Comparación de métodos por índice de Rand ajustado.}
\end{figure}

\begin{figure}[H]
\centering
\includegraphics[width=0.75\linewidth]{\detokenize{../results/figures/top_marker_gene_heatmap.png}}
\caption{Mapa de calor de genes marcadores por grupo.}
\end{figure}

\section{Discusión}
El resultado apoya la hipótesis, pero con cuidado. Los transcriptomas sanguíneos
sí contienen información relacionada con severidad, aunque no separan a todos los
pacientes sin error. Esto tiene sentido porque UCI/no UCI es una etiqueta clínica
práctica, no un estado molecular puro. La parte más fuerte del análisis es la
coherencia entre el grupo con más pacientes UCI, los genes marcadores y los
términos funcionales de inflamación/respuesta innata.

Limitaciones principales:
\begin{itemize}
\item Se usaron matrices procesadas de GEO, no lecturas crudas.
\item No se ajustó por edad, sexo, ventilación mecánica ni comorbilidades.
\item Los marcadores son exploratorios y no sustituyen DESeq2 o limma-voom.
\item Sangre completa mezcla composición celular y activación transcripcional.
\item g:Profiler depende de bases de datos externas que pueden cambiar.
\end{itemize}

\section{Conclusión}
El agrupamiento no supervisado recuperó una señal parcial de UCI/no UCI en
GSE157103. El mejor resultado fue agrupamiento jerárquico
con ward, k=2, ARI=0.263 y exactitud binaria
aproximada de 0.76. El resultado no funciona como
diagnóstico perfecto, pero sí muestra una señal biológica consistente:
inflamación, respuesta innata y degranulación de neutrófilos aparecen en el
grupo con mayor proporción de pacientes UCI.

\section*{Bibliografía}
\begin{enumerate}
\item Overmyer, K. A. et al. (2021). Large-Scale Multi-omic Analysis of COVID-19 Severity. \textit{Cell Systems}, 12(1), 23--40.e7.
\item NCBI GEO. GSE157103: Large-scale Multi-omic Analysis of COVID-19 Severity.
\item Pedregosa, F. et al. (2011). Scikit-learn: Machine Learning in Python. \textit{Journal of Machine Learning Research}, 12, 2825--2830.
\item McInnes, L. et al. (2018). UMAP: Uniform Manifold Approximation and Projection. \textit{Journal of Open Source Software}, 3(29), 861.
\item Raudvere, U. et al. (2019). g:Profiler: a web server for functional enrichment analysis. \textit{Nucleic Acids Research}, 47(W1), W191--W198.
\item Benjamini, Y. y Hochberg, Y. (1995). Controlling the False Discovery Rate. \textit{Journal of the Royal Statistical Society: Series B}, 57(1), 289--300.
\end{enumerate}

\end{document}
